\chapter{Reference frames and planar rotations}
\label{chap:lecture02}

A robot must relate what its sensors observe to a description of the surrounding
world. An obstacle may be directly ahead of the robot, yet lie in a different
direction when described on a map. Neither description is ambiguous once the
reference frame is specified. The task is to express the same geometric point
in different coordinate systems.

This chapter develops that change of coordinates for planar frames with a
common origin. Their relative orientation is described by one angle, or
equivalently by a rotation matrix. The matrix formulation makes it possible to
transform coordinates, reverse a transformation, and compose transformations
through several frames using the linear algebra of Chapter~\ref{chap:lecture01}.

\section{Coordinates depend on a reference frame}
A planar Cartesian reference frame consists of an origin and two perpendicular
unit directions, denoted $\boldsymbol{x}$ and $\boldsymbol{y}$. A world frame
$W$ provides a fixed reference for a map, while a robot frame $R$ is attached
to the robot. A sensor can have its own frame $S$.

\begin{figure}[htbp]
\centering\begin{tikzpicture}[scale=0.85]
\draw[rounded corners,gray] (-.5,-.5) rectangle (7,3.6);
\foreach \x/\y in {2.8/2.6,4.4/2.6,6/2.6,3.2/1.2,4.8/1.2,6.2/.4}
  \draw[fill=gray!15] (\x,\y) circle (.23);
\draw[axis] (0,0)--(1.6,0) node[right] {$x_W$};
\draw[axis] (0,0)--(0,1.6) node[above] {$y_W$};
\node[below left] at (0,0) {$O_W$};
\begin{scope}[shift={(1.6,1.4)},rotate=30]
\draw[fill=robotframe!10,rounded corners] (-.35,-.25) rectangle (.35,.25);
\draw[axis,robotframe] (0,0)--(1.2,0) node[right] {$x_R$};
\draw[axis,robotframe] (0,0)--(0,1.2) node[above] {$y_R$};
\end{scope}
\node[below] at (1.6,1.05) {robot};
\node at (4.5,3.15) {obstacles in a world map};
\end{tikzpicture}

\caption{A map and a robot use different reference frames. In a general setting
their origins and orientations both differ.}
\label{fig:l02-map}
\end{figure}

If the robot's $x_R$ axis points forward and $y_R$ points left, an obstacle one
unit ahead has robot coordinates $[1\;0]\trans$. An obstacle one unit to the
left has coordinates $[0\;1]\trans$. These descriptions alone do not specify
where the obstacles belong on the world map; the relation between the frames
is also needed.

\subsection{Notation and the common-origin assumption}
Let $B$ denote a geometric point. Its coordinates are column vectors,
\begin{equation}
 B^W=\begin{bmatrix}x_B^W\\y_B^W\end{bmatrix},\qquad
 B^R=\begin{bmatrix}x_B^R\\y_B^R\end{bmatrix}.
\end{equation}
The superscript states the frame in which the coordinates are expressed.
The point $B$ does not change when its coordinate description changes.

For now, assume that $W$ and $R$ share an origin $O$ and have the same
orientation convention for their axes. Let $\theta=\theta_R^W$ be the signed
angle from $x_W$ to $x_R$, positive counterclockwise. Only a rotation separates
the two frames. Figure~\ref{fig:l02-coordinates} shows how projections onto
their axes produce different coordinates for the same point.

\begin{figure}[htbp]
\centering\begin{tikzpicture}[scale=1.05]
\coordinate (B) at (1.75,2.35);
\draw[axis] (0,0)--(3.8,0) node[right] {$x_W$};
\draw[axis] (0,0)--(0,3.25) node[above] {$y_W$};
\draw[axis,robotframe] (0,0)--(30:3.8) node[right] {$x_R$};
\draw[axis,robotframe] (0,0)--(120:2.6) node[above] {$y_R$};
\draw[projection] (B)--(1.75,0) node[below] {$x_B^W$};
\draw[projection] (B)--(0,2.35) node[left] {$y_B^W$};
\coordinate (X) at ({2.69054*cos(30)},{2.69054*sin(30)});
\coordinate (Y) at ({1.16016*cos(120)},{1.16016*sin(120)});
\draw[projection,robotframe] (B)--(X) node[below right] {$x_B^R$};
\draw[projection,robotframe] (B)--(Y) node[left] {$y_B^R$};
\draw[thick,sensorframe] (0,0)--(B);
\fill[sensorframe] (B) circle (2pt) node[above right] {$B$};
\draw[-{Stealth}] (.8,0) arc (0:30:.8);
\node at (15:1.02) {$\theta$};
\node[below left] at (0,0) {$O$};
\end{tikzpicture}

\caption{A point expressed in two frames with a common origin. Black axes
belong to $W$ and red axes to $R$; dashed lines give the coordinate projections.}
\label{fig:l02-coordinates}
\end{figure}

\section{Constructing the rotation matrix}
\subsection{The columns are the rotated basis vectors}
The direction $\boldsymbol{x}_R$ makes angle $\theta$ with $x_W$. Because it
is a unit vector, its components in the world frame are
\begin{equation}
 [\boldsymbol{x}_R]^W=\begin{bmatrix}\cos\theta\\\sin\theta\end{bmatrix}.
\end{equation}
The direction $\boldsymbol{y}_R$ is a further $\pi/2$ counterclockwise, so
\begin{equation}
 [\boldsymbol{y}_R]^W
 =\begin{bmatrix}\cos(\theta+\pi/2)\\\sin(\theta+\pi/2)\end{bmatrix}
 =\begin{bmatrix}-\sin\theta\\\cos\theta\end{bmatrix}.
\end{equation}
Here brackets with a superscript mean the components of a geometric vector
in the indicated frame. The negative sign in the first component of
$[\boldsymbol{y}_R]^W$ is essential: for a small positive angle, $y_R$
points to the left of $y_W$.

\begin{figure}[htbp]
\centering\begin{tikzpicture}[scale=2.1]
\draw[axis] (-.85,0)--(1.3,0) node[right] {$x_W$};
\draw[axis] (0,-.15)--(0,1.25) node[above] {$y_W$};
\draw[axis,robotframe] (0,0)--(30:1) node[above right] {$\boldsymbol{x}_R$};
\draw[axis,robotframe] (0,0)--(120:1) node[above left] {$\boldsymbol{y}_R$};
\draw[projection] (30:1)--({cos(30)},0) node[below] {$\cos\theta$};
\draw[projection] (30:1)--(0,.5) node[right,fill=white,inner sep=1pt] {$\sin\theta$};
\draw[projection] (120:1)--(-.5,0) node[below] {$-\sin\theta$};
\draw[projection] (120:1)--(0,{sin(120)}) node[right] {$\cos\theta$};
\draw[-{Stealth}] (.32,0) arc (0:30:.32);
\node at (15:.48) {$\theta$};
\begin{scope}[shift={(2,0.15)},scale=.8]
\draw[thick] (0,0)--(1.1,0)--(1.1,.635)--cycle;
\node[below] at (.55,0) {$\cos\theta$};
\node[right] at (1.1,.32) {$\sin\theta$};
\node[above left] at (.55,.32) {$1$};
\draw (.35,0) arc (0:30:.35);
\node at (15:.55) {$\theta$};
\draw (.99,0)--(.99,.11)--(1.1,.11);
\end{scope}
\end{tikzpicture}

\caption{Components of the robot's unit basis vectors in the world frame.
The right triangle explains the sine and cosine components of $\boldsymbol{x}_R$.}
\label{fig:l02-basis}
\end{figure}

Putting these two column vectors side by side defines the rotation matrix:
\begin{equation}
 \boxed{R_R^W
 =\begin{bmatrix}[\boldsymbol{x}_R]^W &[\boldsymbol{y}_R]^W\end{bmatrix}
 =\begin{bmatrix}\cos\theta&-\sin\theta\\\sin\theta&\cos\theta\end{bmatrix}.}
 \label{eq:l02-rotation}
\end{equation}
The subscript identifies the frame whose orientation is being described, and
the superscript identifies the reference frame. As a coordinate map, $R_R^W$
takes components expressed in $R$ and returns components expressed in $W$.
An angle is a compact description of orientation; the matrix provides a linear
map on coordinate vectors once that angle is fixed.

\subsection{Orthogonality}
Write $c=\cos\theta$ and $s=\sin\theta$. Direct multiplication gives
\[
 (R_R^W)\trans R_R^W
 =\begin{bmatrix}c&s\\-s&c\end{bmatrix}
  \begin{bmatrix}c&-s\\s&c\end{bmatrix}
 =\begin{bmatrix}c^2+s^2&0\\0&c^2+s^2\end{bmatrix}=I_2.
\]
Thus $R_R^W$ is orthogonal and
\begin{equation}
 (R_R^W)^{-1}=(R_R^W)\trans
 =\begin{bmatrix}\cos\theta&\sin\theta\\-\sin\theta&\cos\theta\end{bmatrix}.
 \label{eq:l02-inverse}
\end{equation}
Its determinant is $c^2+s^2=1$. This last property distinguishes a rotation
from an orthogonal reflection, whose determinant is $-1$.

\section{Transforming the coordinates of a point}
The vector from the common origin to $B$ can be decomposed along the robot axes:
\[
 \overrightarrow{OB}=x_B^R\boldsymbol{x}_R+y_B^R\boldsymbol{y}_R.
\]
Expressing both basis vectors in $W$ gives
\begin{equation}
 \boxed{B^W=R_R^W B^R.}
 \label{eq:l02-coordinate-map}
\end{equation}
In scalar form, this is
\begin{align}
 x_B^W&=\cos\theta\,x_B^R-\sin\theta\,y_B^R,\\
 y_B^W&=\sin\theta\,x_B^R+\cos\theta\,y_B^R.
\end{align}
The dimensions are $(2\times2)(2\times1)=(2\times1)$, and the input frame
$R$ agrees with the matrix subscript. Both checks are useful: correct matrix
dimensions alone do not establish that the intended frames match.

\subsection{A geometric derivation}
For $B\ne O$, let $d$ be the distance from $O$ to $B$, let $\alpha$ be the
angle from $x_R$ to $\overrightarrow{OB}$, and let $\beta$ be the angle from
$x_W$ to that vector. The angles satisfy $\beta=\alpha+\theta$ (modulo $2\pi$).
Consequently,
\[
 x_B^R=d\cos\alpha,\quad y_B^R=d\sin\alpha,\qquad
 x_B^W=d\cos\beta,\quad y_B^W=d\sin\beta.
\]

\begin{figure}[htbp]
\centering\begin{tikzpicture}[scale=.95]
\coordinate (B) at (58:3.6);
\draw[axis] (0,0)--(4,0) node[right] {$x_W$};
\draw[axis] (0,0)--(0,3.65) node[above] {$y_W$};
\draw[axis,robotframe] (0,0)--(25:4) node[right] {$x_R$};
\draw[axis,robotframe] (0,0)--(115:2.8) node[above] {$y_R$};
\draw[thick,sensorframe] (0,0)--(B) node[midway,above left] {$d$};
\fill (B) circle (2pt) node[above right] {$B$};
\draw[projection] (B)--({3.6*cos(58)},0) node[below] {$x_B^W$};
\draw[projection] (B)--(0,{3.6*sin(58)}) node[left] {$y_B^W$};
\draw[-{Stealth}] (.8,0) arc (0:25:.8);
\node at (12.5:1.05) {$\theta$};
\draw[-{Stealth},robotframe] (25:1.3) arc (25:58:1.3);
\node[robotframe] at (41.5:1.57) {$\alpha$};
\draw[-{Stealth},sensorframe] (2.35,0) arc (0:58:2.35);
\node[sensorframe,fill=white,inner sep=2pt] at (40:2.35) {$\beta$};
\node at (3.45,3.2) {$\beta=\alpha+\theta$};
\end{tikzpicture}

\caption{The same position vector has angle $\alpha$ in $R$ and angle
$\beta=\alpha+\theta$ in $W$. Its length $d$ is independent of the frame.}
\label{fig:l02-angles}
\end{figure}

Expanding the angle sums gives
\begin{align*}
 x_B^W&=d\cos(\alpha+\theta)
 =d\cos\alpha\cos\theta-d\sin\alpha\sin\theta\\
 &=\cos\theta\,x_B^R-\sin\theta\,y_B^R,\\
 y_B^W&=d\sin(\alpha+\theta)
 =d\cos\alpha\sin\theta+d\sin\alpha\cos\theta\\
 &=\sin\theta\,x_B^R+\cos\theta\,y_B^R.
\end{align*}
These are precisely the entries of~\eqref{eq:l02-coordinate-map}. The matrix
formula also applies at $B=O$, where both coordinate vectors are zero and a
direction angle is unnecessary.

\subsection{Worked example: a thirty-degree orientation}
Suppose
\[
 B^R=\begin{bmatrix}0.8\\0.5\end{bmatrix},\qquad \theta=\frac{\pi}{6}.
\]
Then
\begin{align}
 B^W
 &=\begin{bmatrix}\sqrt3/2&-1/2\\1/2&\sqrt3/2\end{bmatrix}
   \begin{bmatrix}0.8\\0.5\end{bmatrix}\\
 &=\begin{bmatrix}0.4\sqrt3-0.25\\0.4+0.25\sqrt3\end{bmatrix}
 \approx\begin{bmatrix}0.4428\\0.8330\end{bmatrix}.
 \label{eq:l02-example}
\end{align}
The coordinates change, but the point's distance from the common origin does
not. Indeed, orthogonality implies
\[
 (B^W)\trans B^W=(B^R)\trans(R_R^W)\trans R_R^W B^R
 =(B^R)\trans B^R=0.8^2+0.5^2=0.89.
\]
This is a change of coordinate description of a fixed point, rather than a
physical motion of that point.

\section{Reversing a coordinate transformation}
To recover robot coordinates from world coordinates, multiply
Equation~\eqref{eq:l02-coordinate-map} on the left by the inverse:
\[
 (R_R^W)^{-1}B^W=(R_R^W)^{-1}R_R^W B^R=B^R.
\]
Using orthogonality,
\begin{equation}
 \boxed{B^R=(R_R^W)\trans B^W=R_W^R B^W,\qquad
 R_W^R=(R_R^W)\trans.}
\end{equation}
Reversing the frame relation exchanges the frame indices and negates the
relative angle. If $R(\theta)$ denotes the matrix in
Equation~\eqref{eq:l02-rotation}, then $R(\theta)\trans=R(-\theta)$.
For arbitrary frames $A$ and $F$, the same rule reads
\begin{equation}
 R_F^A=(R_A^F)\trans.
 \label{eq:l02-swap}
\end{equation}
For example, applying the transpose to the exact vector in
Equation~\eqref{eq:l02-example} recovers $[0.8\;0.5]\trans$.

\section{Composing rotations through several frames}
A measurement may begin in a sensor frame $S$, pass through a robot frame $R$,
and finally be expressed in $W$. Assume again that all origins coincide.
Let the orientation of $R$ relative to $W$ be $\theta$, and the orientation
of $S$ relative to $R$ be $\varphi$.

\begin{figure}[htbp]
\centering\begin{tikzpicture}[scale=1.05]
\draw[axis] (0,0)--(3.7,0) node[right] {$x_W$};
\draw[axis] (0,0)--(0,3.4) node[above] {$y_W$};
\draw[axis,robotframe] (0,0)--(25:3.7) node[right] {$x_R$};
\draw[axis,robotframe] (0,0)--(115:3) node[above] {$y_R$};
\draw[axis,sensorframe] (0,0)--(55:3.7) node[above right] {$x_S$};
\draw[axis,sensorframe] (0,0)--(145:3) node[above left] {$y_S$};
\draw[-{Stealth},robotframe] (.8,0) arc (0:25:.8);
\node[robotframe] at (12.5:1.05) {$\theta$};
\draw[-{Stealth},sensorframe] (25:1.5) arc (25:55:1.5);
\node[sensorframe] at (40:1.76) {$\varphi$};
\draw[thick] (0,0)--(75:3.2);
\fill (75:3.2) circle (2pt) node[right] {$B$};
\node[below left] at (0,0) {$O$};
\end{tikzpicture}

\caption{Three frames with a common origin. The sensor frame (blue) is rotated
by $\varphi$ relative to the robot frame (red), which is rotated by $\theta$
relative to the world frame (black).}
\label{fig:l02-composition}
\end{figure}

The two successive coordinate transformations are
\[
 B^R=R_S^R B^S,\qquad B^W=R_R^W B^R.
\]
Substitution gives
\begin{equation}
 B^W=R_R^W R_S^R B^S,\qquad
 \boxed{R_S^W=R_R^W R_S^R.}
 \label{eq:l02-chain}
\end{equation}
The rightmost matrix acts first. Its output is expressed in $R$, precisely the
frame expected by the matrix to its left. In two dimensions,
\[
 R_R^W=R(\theta),\qquad R_S^R=R(\varphi),\qquad
 R_S^W=R(\theta+\varphi),
\]
as can be verified with the sine and cosine addition identities.

\subsection{Reading the frame indices}
The general composition rule is
\begin{equation}
 R_C^A=R_B^A R_C^B.
\end{equation}
The subscript of the left factor matches the superscript of the next factor.
The remaining indices identify the overall output and input frames. This is
a rule about the meanings of the maps, not ordinary cancellation of numbers.

For example,
\begin{equation}
 R_2^3 R_4^2 R_5^4 R_1^5=R_1^3.
 \label{eq:l02-long-chain}
\end{equation}
Reading from right to left, the coordinates pass through frames
$1\to5\to4\to2\to3$. By contrast, in the written product $R_2^3 R_3^4$,
the inner frame labels $2$ and $4$ do not match. Its factors cannot be composed
as that sequence of coordinate maps merely by cancelling the repeated $3$.

Planar rotation matrices happen to commute numerically:
$R(\theta)R(\varphi)=R(\varphi)R(\theta)$. This special property does not
change the input and output frame labels of an individual map. Keeping the
frame chain explicit also prepares the notation for three-dimensional
rotations, where commutativity does not hold in general.

\subsection{Simplifying chains containing transposes}
First apply the transpose rules, then inspect the frame chain. Consider
\[
 R_2^3 (R_2^4)\trans (R_5^1 R_4^5)\trans.
\]
Transposing the last product reverses its factor order, so
\begin{align*}
 R_2^3 (R_2^4)\trans (R_5^1 R_4^5)\trans
 &=R_2^3 R_4^2 (R_4^5)\trans (R_5^1)\trans\\
 &=R_2^3 R_4^2 R_5^4 R_1^5\\
 &=R_1^3.
\end{align*}
Two different operations are involved: a transpose exchanges the frame indices
of one rotation matrix, and transposing a product reverses the order of its
factors. Both are necessary to obtain a consistent chain.

\section{The limit of a rotation-only model}
The formula $B^W=R_R^W B^R$ describes point coordinates when the origins
coincide. In a general robot configuration, the origin $O_R$ is displaced from
$O_W$, as in Figure~\ref{fig:l02-translation}. A rotation matrix describes the
relative orientation, but contains no information about this displacement.

\begin{figure}[htbp]
\centering\begin{tikzpicture}[scale=.9]
\draw[axis] (0,0)--(5,0) node[right] {$x_W$};
\draw[axis] (0,0)--(0,3.3) node[above] {$y_W$};
\node[below left] at (0,0) {$O_W$};
\draw[axis,sensorframe] (0,0)--(2.6,1.5) node[midway,above left] {$\boldsymbol{t}_R^W$};
\begin{scope}[shift={(2.6,1.5)}]
\fill (0,0) circle (2pt) node[below right] {$O_R$};
\draw[axis,robotframe] (0,0)--(30:2) node[right] {$x_R$};
\draw[axis,robotframe] (0,0)--(120:1.8) node[above] {$y_R$};
\end{scope}
\end{tikzpicture}

\caption{Different origins introduce a translation $\boldsymbol{t}_R^W$.
Orientation alone cannot determine the world coordinates of a point measured
relative to $O_R$.}
\label{fig:l02-translation}
\end{figure}

To see the limitation, take the point $B=O_R$. Its robot coordinates are zero,
so $R_R^W B^R=0$ for every rotation matrix. Its world coordinates, however,
are the nonzero displacement of $O_R$ from $O_W$. A complete point-coordinate
transformation must therefore include translation as well as rotation.
The rotation rules remain applicable to direction vectors, which do not depend
on the choice of origin. Distinguishing point positions from directions is
essential when extending the model.

\section{Key identities}
\begin{center}
\begin{tabular}{ll}
\toprule
Operation & Identity \\
\midrule
Planar rotation & $R(\theta)=\begin{bmatrix}\cos\theta&-\sin\theta\\\sin\theta&\cos\theta\end{bmatrix}$ \\[12pt]
Point coordinates, common origin & $B^W=R_R^W B^R$ \\
Reverse transformation & $R_W^R=(R_R^W)\trans=R(-\theta)$ \\
Composition & $R_S^W=R_R^W R_S^R$ \\
Planar angle composition & $R(\theta)R(\varphi)=R(\theta+\varphi)$ \\
\bottomrule
\end{tabular}
\end{center}

\clearpage
\section{Additional exercises}
Use column vectors, show your calculations, and label the input and output
frames. Unless stated otherwise, all frames share an origin.
\begin{enumerate}
\item \textbf{Sketching frames and coordinates.}
Draw world axes $x_W,y_W$ and robot axes rotated counterclockwise by $\pi/2$.
A point has $B^R=[2\;1]\trans$. Locate it on your sketch, compute $B^W$,
and check that your answer agrees with its position relative to the world axes.

\item \textbf{Constructing a rotation matrix.}
The robot frame is rotated by $-\pi/4$ relative to the world frame. Write
$[\boldsymbol{x}_R]^W$ and $[\boldsymbol{y}_R]^W$, and use them to construct
$R_R^W$. Verify $(R_R^W)\trans R_R^W=I_2$ and $\det(R_R^W)=1$.

\item \textbf{Forward and reverse transformations.}
Let $\theta_R^W=\pi/6$ and $B^R=[1\;2]\trans$. Find $B^W$ in exact form.
Construct $R_W^R$ and use it to recover $B^R$. Verify that the squared distance
from the common origin is the same in both coordinate descriptions.

\item \textbf{From a sensor to the world.}
Suppose $\theta_R^W=\pi/6$, $\theta_S^R=\pi/3$, and
$B^S=[2\;-1]\trans$. Compute $R_S^W$ and $B^W$. Check your result by first
finding $B^R$ and then transforming it to $W$. What is the orientation of $S$
relative to $W$?

\item \textbf{Frame chains and transposes.}
Simplify the first two expressions to a single rotation matrix. For the third,
explain why the written frame chain is inconsistent, even though the numerical
matrix dimensions permit multiplication:
\[
 R_B^A R_C^B R_D^C,\qquad
 (R_A^B)\trans (R_C^D R_B^C)\trans,\qquad
 R_B^A R_D^C.
\]
Assume distinct letters denote distinct frames. State the input and output
frames of each valid chain.

\item \textbf{Why an origin matters.}
The robot axes are parallel to the world axes, but the robot origin has world
coordinates $[3\;1]\trans$. A point has robot coordinates $[2\;0]\trans$.
Use a sketch to find its world coordinates. What would the rotation-only
formula predict, and what information would it omit?

\item \textbf{Length preservation.}
Let $v^W=R_R^W v^R$ be two coordinate descriptions of a direction vector.
Using orthogonality, show that $(v^W)\trans v^W=(v^R)\trans v^R$.
Explain geometrically why a change of reference frame cannot change the
length of the vector.
\end{enumerate}
